\documentclass[11pt]{article}
\usepackage[a4paper,margin=1in]{geometry}
\usepackage{graphicx}
\usepackage{listings}
\usepackage{xcolor}
\usepackage{booktabs}
\usepackage{array}
\usepackage{hyperref}
\usepackage{parskip}
\hypersetup{colorlinks=true, linkcolor=black, urlcolor=blue}

\definecolor{codebg}{RGB}{245,245,245}
\lstdefinestyle{cstyle}{
  language=C,
  backgroundcolor=\color{codebg},
  basicstyle=\ttfamily\footnotesize,
  keywordstyle=\color{blue}\bfseries,
  commentstyle=\color{gray}\itshape,
  stringstyle=\color{teal},
  numbers=left,
  numberstyle=\tiny\color{gray},
  stepnumber=1,
  breaklines=true,
  frame=single,
  tabsize=2,
  showstringspaces=false
}
\lstdefinestyle{outstyle}{
  backgroundcolor=\color{codebg},
  basicstyle=\ttfamily\footnotesize,
  breaklines=true,
  frame=single,
  showstringspaces=false
}

\title{\vspace{-2cm}}
\date{}

\begin{document}
\begin{titlepage}
\centering
\vspace*{1cm}
{\large \textbf{Sri Sivasubramaniya Nadar College of Engineering}}\\
{\small (An Autonomous Institution, Affiliated to Anna University, Chennai)}\\
{\small Kalavakkam -- 603 110}\\[0.3cm]
{\small Department of Computer Science \& Engineering}\\[1.5cm]

\rule{\linewidth}{0.5mm}\\[0.4cm]
{\LARGE \textbf{Computer Networks Laboratory}}\\[0.3cm]
{\Large 1513 -- Computer Networks Lab}\\[0.4cm]
\rule{\linewidth}{0.5mm}\\[1.5cm]

{\huge \textbf{Assignment 6}}\\[0.4cm]
{\Large Implementation of Error Detection using Checksum and CRC}\\[2cm]

\begin{center}
\begin{tabular}{ll}
\textbf{Name:} & Sundareswaran Rajaguru \\[0.2cm]
\textbf{Register Number:} & 3122247001048 \\[0.2cm]
\textbf{Year / Semester:} & III Year / V Semester \\[0.2cm]
\textbf{Course:} & 1513 Computer Networks Lab \\[0.2cm]
\textbf{Date of Submission:} & 19 August 2026 \\[0.2cm]
\end{tabular}
\end{center}

\vfill
\end{titlepage}

\section*{1. Objective}
To develop a socket programming application in C that simulates reliable data transmission using
two error-detection techniques:
\begin{itemize}
  \item \textbf{Checksum:} data is divided into 8-bit words, summed using 1's complement addition
        with end-around carry, and the resulting checksum is verified at the receiver.
  \item \textbf{Cyclic Redundancy Check (CRC):} data is treated as a binary polynomial, divided by
        a generator polynomial $G(x) = x^4+x+1$ (i.e.\ \texttt{10011}), and the remainder (CRC) is
        appended to the frame before transmission over TCP sockets. Bit errors are then
        deliberately introduced to demonstrate that CRC verification correctly detects them.
\end{itemize}

\section*{2. Server Algorithm}
\begin{enumerate}
  \item Create a TCP socket, bind to port 8080, and \texttt{listen()} for connections.
  \item \texttt{accept()} each client connection and spawn a detached thread to service it
        concurrently.
  \item Read one byte identifying the mode: \texttt{'C'} for Checksum or \texttt{'R'} for CRC.
  \item \texttt{recv()} the payload, where the final byte is the checksum/CRC appended by the
        sender and the preceding bytes are the data.
  \item \textbf{Checksum mode:} recompute the 1's-complement checksum of the received data
        (\texttt{checksum\_verify()}) and compare it against the received checksum byte.
  \item \textbf{CRC mode:} recompute the CRC of the received data using the same generator
        polynomial \texttt{10011} (\texttt{crc\_calc()}) and compare it against the received CRC
        byte.
  \item Print the received vs.\ calculated value and declare the frame \texttt{ERROR-FREE} if they
        match, or \texttt{ERROR DETECTED} otherwise.
\end{enumerate}

\section*{3. Client Algorithm}
\begin{enumerate}
  \item Ask the user to choose Checksum or CRC mode, then read the data string to transmit.
  \item \textbf{Checksum mode:} sum each byte into a running 16-bit accumulator, printing the
        intermediate sum after every byte; fold the carry out of the top byte back into the sum
        (end-around carry); take the 1's complement of the result as the final checksum.
  \item \textbf{CRC mode:} process the data one byte at a time with an 8-bit polynomial-division
        register: XOR each byte into the register, then shift left 8 times, XOR-ing in the
        generator polynomial (\texttt{0x13} = \texttt{10011}) whenever the top bit is 1; the
        register's final value is the CRC remainder.
  \item Connect to the server over TCP, send the mode byte, then send the data with the
        checksum/CRC appended as the last byte.
  \item \textbf{CRC mode only:} optionally flip one bit of the already-sent frame at a
        user-chosen byte position and resend it to the server unmodified except for that
        corruption, to demonstrate that the receiver's recomputed CRC no longer matches.
\end{enumerate}

\section*{4. Server Code}
\subsection*{4.1 server.c}
\lstinputlisting[style=cstyle]{server.c}

\section*{5. Client Code}
\subsection*{5.1 client.c}
\lstinputlisting[style=cstyle]{client.c}

\section*{6. Output / Test Cases}
Three cases were exercised: (1) Checksum on the string \texttt{HELLO}, verified error-free; (2)
CRC on the string \texttt{TEST}, verified error-free; (3) the same CRC frame re-sent with one bit
of the first byte flipped, which the server correctly flags as \texttt{ERROR DETECTED}.

\subsection*{6.1 Client Output -- Checksum Mode (HELLO)}
\lstinputlisting[style=outstyle]{output_checksum.txt}

\subsection*{6.2 Client Output -- CRC Mode with Injected Bit Error (TEST)}
\lstinputlisting[style=outstyle]{output_crc.txt}

\subsection*{6.3 Server-side Verification Log}
\begin{lstlisting}[style=outstyle]
[Server] Listening on port 8080...

[Server] Received: HELLO (6 bytes)
[CHECKSUM] Received: 0x8A, Calculated: 0x8A
[Result] ERROR-FREE

[Server] Received: TEST (5 bytes)
[CRC] Received: 0xF1, Calculated: 0xF1
[Result] ERROR-FREE

[Server] Received: UEST (5 bytes)
[CRC] Received: 0xF1, Calculated: 0xF6
[Result] ERROR DETECTED
\end{lstlisting}
Flipping bit 0 of the character \texttt{'T'} (\texttt{0x54}) changed it to \texttt{'U'}
(\texttt{0x55}), which the receiver's Wireshark-independent, application-level CRC check
correctly identified: the received CRC (\texttt{0xF1}, computed by the sender over the original
data) no longer matches the CRC the receiver recalculates over the corrupted data
(\texttt{0xF6}).

\section*{7. Learning Outcomes}
\begin{itemize}
  \item Implemented 1's-complement checksum arithmetic with end-around carry from first
        principles, and observed step-by-step how the carry-out of the high byte must be folded
        back into the sum before complementing.
  \item Implemented binary polynomial division (CRC) using an 8-bit shift-and-XOR register driven
        by the generator polynomial $G(x)=x^4+x+1$, and understood why the choice of generator
        determines which burst-error patterns are detectable.
  \item Verified experimentally that CRC reliably detects single-bit corruption, whereas a naive
        recomputation and comparison is only meaningful if both sender and receiver use the exact
        same generator polynomial and byte framing.
  \item Practised structuring an error-detection scheme as an application-layer protocol over TCP,
        with the detection code appended to the payload as a trailer, mirroring how data-link
        layer frames carry an FCS/CRC field.
  \item Gained experience building a multi-mode (checksum/CRC) client-server system that shares a
        single TCP port and thread-per-client server design.
\end{itemize}

\end{document}
